\summarychapter{34}{A Brand from the Burning}{Wesley and Methodism}

\begin{inbrief}
In the Age of Reason a spiritual renewal, the Evangelical Awakening, changed three regions: Germany
through Pietism, the British Isles through the Methodists and the American colonies through the Great
Awakening. The chapter follows John Wesley from his strict upbringing and his failure in Georgia to
his heart-warming at Aldersgate, field preaching, the organization of the Methodist societies and
their separation into a new denomination.
\question{Who were the Evangelicals, and how did they gain major significance in Christian history?}
\end{inbrief}

Late in January 1736 the ship \emph{Simmonds}, bound for Savannah, ran into violent Atlantic storms.
A slight young Anglican clergyman on board, John Wesley, was frozen with fear; he had preached
salvation to others but was afraid to die. He was awed by a company of Moravians from Herrnhut who,
as the sea split the mainsail, calmly went on singing psalms. Asked afterward whether they had been
frightened, one answered that even their women and children were not afraid to die. Wesley called it
the most glorious day he had seen. He had the form of godliness but not yet its power.

\ssection{The Evangelical Awakening}
The Methodist revival explains the origin of a denomination that now numbers over twenty million and
also throws light on evangelical Christianity as a whole. The singing Moravians represent its Pietist
wing; the main body came from Britain and its colonies. In the 1730s a sudden concern to preach to
the unconverted broke out in several places at once. Jonathan Edwards in Northampton, Massachusetts,
Ebenezer and Ralph Erskine in Scotland, Howel Harris in Wales and George Whitefield in England all
came before Wesley. Most evangelical beliefs were Puritan: human sinfulness, Christ's atoning death,
unmerited grace, the salvation of the true believer. But the Puritans were more concerned with
politics and the holy commonwealth. The Evangelicals, less detached from politics than the Pietists,
had one controlling passion, the conversion of the lost, and cared less for reforming churches than
for preaching to everyone: nominal Christians, scoffers and heathen.

England was an unlikely place for it. Among the rich and educated the Enlightenment had pushed
religion to the edge of life, and in the Anglican church as in the Nonconformist bodies, Baptist and
Congregationalist, Puritan zeal seemed a thing of the past. The country had had its fill of holy
causes and wanted moderation. Voltaire described an English sermon as a solid, sometimes dry
dissertation read without gesture. Ministers passed over human sinfulness and approached God with
gentle awe and cheerfulness, as in Addison's hymn ``The Spacious Firmament on High.'' The
Latitudinarians, a large party in the Church of England that included John Tillotson, archbishop of
Canterbury (1691--1694), denounced ``enthusiasm,'' meaning any emotion stirred by fervent preaching,
and stressed proper conduct: be generous, humane and tolerant, and avoid bigotry and fanaticism.

\ssection{The Little Giant}
John Wesley (1703--1791) came from a home of decency and order that joined Anglican and Nonconformist
piety. His father Samuel was a learned high-church rector at Epworth in Lincolnshire. His mother
Susanna, daughter of a Nonconformist minister, bore nineteen children, of whom John was the
fifteenth; she taught them ``to fear the rod and cry softly'' and gave each child separate religious
instruction every week, which required method. When he was six the rectory burned, and he was rescued
from an upper window by a neighbor standing on another's shoulders. Ever after he called himself ``a
brand plucked from the burning'' and never doubted God's hand on his life.

At seventeen he went to Oxford, to Christ Church and later Lincoln College. He read widely. The early
Greek fathers taught him that the goal of the Christian life is perfection, understood as a process
of disciplined love rather than a state; Jeremy Taylor, Thomas \`a Kempis and William Law taught him
that it is the consecration of the whole person in love to God and neighbor. They convinced him ``of
the absolute impossibility of being half a Christian,'' and he listed his weaknesses and made rules
against them. In 1726 he was elected a fellow of Lincoln, which gave him standing and an income; two
years later he was ordained and for a time assisted his father. Back at Oxford he found that his
brother Charles, alarmed by deism, had gathered a few students who meant to take religion seriously.
John became their leader, and they adopted a rule of prayer, Bible reading and frequent Communion.
Other undergraduates jeered at the Holy Club, the Bible moths, the Methodists; the last name stuck.
They gave to the poor and visited prisoners, and were joined by the bright and brash George
Whitefield of Pembroke College. But John admitted that he lacked the inward peace of a true
Christian.

Then a friend, Dr.~John Burton, suggested Georgia, General James Oglethorpe's new colony, with Charles
as the general's secretary and John as chaplain. Hoping to preach to the Indians, the brothers sailed
in October. The venture was a fiasco. John found the noble savages ``gluttons, thieves, liars and
murderers,'' and the colonists resented his rigid high-church ways, his refusal to bury a
Nonconformist and his ban on fine dresses and gold jewelry in church. An unhappy love affair made
things worse. Sophy Hopkey, the eighteen-year-old niece of Savannah's chief magistrate, tired of his
indecision and eloped with a rival; Wesley barred her from Communion, and her husband sued him for
defamation. After six months of harassment he fled. ``I went to America to convert the Indians,'' he
wrote on the way home, ``but, oh, who shall convert me?''

\ssection{The Holy Heart Warming}
He landed in England on February 1, 1738, discredited and unsure of his faith. Twelve years of
toiling toward perfection had ended in spiritual bankruptcy. The one good thing from Georgia was the
Moravians, and he set out to learn their secret. In London Peter B\"ohler, a young Moravian preacher,
pressed on him his need of a new birth, a personal faith that overcomes sin; justification by faith
is not only a doctrine but an experience of God's forgiveness. On May 24, 1738, Wesley went ``very
unwillingly'' to a society in Aldersgate Street where Luther's preface to Romans was being read.
About a quarter to nine, he wrote, ``I felt my heart strangely warmed. I felt that I did trust in
Christ, Christ alone, for salvation; and an assurance was given me that he had taken away my sins,
even mine.'' His journal shows a continuing insecurity that gave way only when God used his preaching
to great effect, but he had found the assurance and the message that sustained fifty years of
unmatched energy. He still needed a method.

That summer he visited Herrnhut. He met people who showed ``the full assurance of Christian faith,''
but he also noticed self-righteousness and was repelled by the cult around Zinzendorf: ``Isn't the
Count all in all?'' He owed the Moravians the message of justification by faith and the system of
small groups, but he could not be one of them. Back in London his preaching was as zealous and as
fruitless as before. Then, walking from London to Oxford, he read Jonathan Edwards's account of the
conversions at Northampton. It struck him with great force, a direct influence of the American
awakening on the English one. Within weeks came an invitation from George Whitefield. Nine years
younger, Whitefield had followed Wesley to Georgia in 1738 and returned to be ordained. Eager to
reach the masses, he began in February 1739 to preach in the open fields near Bristol to coal miners
who seldom entered a church, and saw ``the white gutters made by their tears'' on their blackened
faces. His preaching was unforgettable. He could make hearers feel the pain of sin and the terror of
hell and then weep with him for forgiveness; the actor David Garrick said he would give a hundred
guineas to say ``Oh!'' like Mr.~Whitefield. He urged Wesley to follow him into the fields. Wesley
hesitated, an Oxford scholar and gentleman so attached to decency and order that he would have
thought ``the saving of souls almost a sin if it had not been done in a church.''

\ssection{To the Fields, to the World}
Against Charles's opposition John went to Bristol, more like a martyr than a glad messenger. He
preached in the open air to more than three thousand, and conversions as real as those in New England
took place on every side. The Methodist revival had begun. The effect on Wesley was as great. B\"ohler
had told him to preach faith until he had it. Aldersgate had taken him from hoping to having, Edwards
and Whitefield had shown him that the Word rightly preached bears visible fruit, and now his own
faith was confirmed by that of his hearers. The anxious, insecure man became a firebrand. That June
he wrote, ``I look upon all the world as my parish.''

He preached in jails, in inns, on ships to Ireland; to thirty thousand in a natural amphitheater in
Cornwall; and, when refused the church at Epworth, to hundreds in the churchyard, standing on his
father's tombstone. He reckoned his travel at no less than 4,500 miles a year, a quarter of a million
in his lifetime, mostly on horseback, reading or preparing sermons with the reins slack. In the early
years he was pelted with stones and sometimes beaten by mobs set on by a hostile squire or parson, but
he feared no one and often awed a crowd by sheer presence. In time the violence stopped, and before
he died china statuettes of him were sold in large numbers. In 1751 he married Molly Vazeille, a
London merchant's widow who had nursed him after a fall on the ice. For two years she tried to travel
with him; then her health and nerves broke, and she left. When she died in 1781 he did not know of it
and was not at the funeral. She had married a man wedded to his mission.

Alone among the leaders of the Awakening, Wesley preached what are now called Arminian beliefs, named
for Jacob Arminius (1560--1609), the Dutch professor who tried to modify Calvinism. He owed Arminius
no special debt, but he firmly opposed predestination, which seemed to him to make God arbitrary. God
wills the salvation of all, and people can accept or refuse his grace. This nearly broke his
friendship with Whitefield, who defended predestination as the guard of God's sovereignty and feared
that Arminianism dulled the sense of sin. They agreed to differ with mutual respect, and at
Whitefield's funeral in 1770 Wesley spoke of his ``most generous and tender friendship.'' But the
Methodists divided into Arminian societies following Wesley and Calvinist ones following Whitefield.

\ssection{The Methodist Structure}
Whitefield had no taste for organization; Wesley was an administrative genius. Methodist societies
sprang up behind him across England, Ireland and Wales. They were not yet congregations. Most members
were Anglicans, and Wesley sent them to their parish churches for worship and Communion. But the
center of their Christian life was the society, where they confessed their sins to one another,
accepted the discipline of their leader, and prayed and sang. Charles Wesley, who had found
forgiveness three days before John, wrote over seven thousand hymns for these meetings, among them
``Jesus, Lover of My Soul''; some think the hymns the revival's greatest legacy.

On the Moravian model John divided the societies into classes of about twelve. The word, from the
Latin \emph{classis}, a division, has nothing to do with school. The classes began as a means of
collecting a penny a week, but he soon saw that the collector could be a spiritual guide and the
members could encourage one another, and the class meeting for testimony and prayer became one of the
most successful features of Methodism. As the work grew he used laymen as preachers and assistants.
He would not call them ministers or let them administer the sacraments; they were his personal
helpers, answerable to him as he was to the Church of England. When by 1744 he could no longer keep
in touch with all of them, he created the Annual Conference, which helped shape policy and doctrine,
always as Wesley decided. He moved his preachers about like a militia and held them to one task,
evangelism and Christian nurture. They were not, he said, the leaders of a sect but messengers to
those who are ``Christians in name but heathens in heart and life.''

By 1748 the ``people called Methodists'' were, like the Pietists, a church within the church. For
forty years Wesley resisted every pressure toward separation: ``I live and die a member of the Church
of England.'' Near the end the needs of America forced his hand. He had sent Francis Asbury to the
colonies, and the first American Methodist conference met in Philadelphia in 1773, but the societies
needed ordained leaders, and the bishop of London would not help. So Wesley appointed two of his lay
preachers, Richard Whatcoat and Thomas Vasey, for America and commissioned Dr.~Thomas Coke as
superintendent, a clear breach of Anglican order. The Christmas Conference at Baltimore in 1784,
which chose Coke and Asbury as superintendents, made the Methodist Church in America a distinct
denomination.

Wesley preached almost to the last and died in London on March 2, 1791, nearly eighty-eight, leaving
seventy-nine thousand followers in England and forty thousand in North America. Afterward the English
Methodists also separated from the Anglican church. But the revival reached far beyond Methodism. It
renewed the religious life of England and its colonies, raised the life of the poor, and stimulated
the overseas missions and the social concern of Evangelicals in the next two centuries.

\begin{keyfacts}
\keypeople{John Wesley; Susanna Wesley; Charles Wesley; George Whitefield; Peter B\"ohler; Jonathan
Edwards; Count von Zinzendorf; John Tillotson; Jacob Arminius; Francis Asbury; Thomas Coke.}
\keydates{%
  \kd{1703}{Birth of John Wesley}
  \kd{1736}{The Wesleys sail to Georgia; the storm and the Moravians}
  \kd{1738}{Wesley's heart ``strangely warmed'' at Aldersgate (May 24)}
  \kd{1739}{Whitefield and Wesley begin field preaching near Bristol}
  \kd{1744}{First Annual Conference of Methodist preachers}
  \kd{1770}{Death of George Whitefield}
  \kd{1773}{First American Methodist conference, Philadelphia}
  \kd{1784}{Christmas Conference at Baltimore: American Methodism a separate church}
  \kd{1791}{Death of John Wesley (March 2)}}
\keyterms{Evangelical Awakening; Evangelicals; Latitudinarians; ``enthusiasm''; Holy Club;
Methodists; new birth; assurance; field preaching; Arminianism; societies; class meeting; lay
preachers; Annual Conference.}
\end{keyfacts}
